% Generated from locale/tr/content/first-order-logic/proof-systems/tableaux.tex; do not hand-edit.


\iftag{FOL}
      {\olfileid[tr]{fol}{prf}{tab}}
      {\olfileid[tr]{pl}{prf}{tab}}

\olsection{\usetoken{P}{tableau}}

Birçok !!{derivation} sistemi !!{sentence}s içeren düzenlemelerle çalışırken,
!!{tableau}s !!{signed formula}s ile çalışır. Bir !!^a{signed formula}, bir
doğruluk değeri işareti ($\True$ veya $\False$) ile !!a{sentence} içeren bir
ikilidir:
\[
\sFmla{\True}{!A} \text{ veya } \sFmla{\False}{!A}.
\]
Bir !!^a{tableau}, aşağı doğru dallanan bir ağaçta düzenlenmiş !!{signed formula}s
içerir. Birkaç \emph{varsayımla} başlar ve üstteki !!{signed formula}s
içinden birine bir çıkarım kuralı uygulanmasıyla elde edilen !!{signed
formula}s ile devam eder. Her kural, !!{signed formula}s arasından birini
veya birkaçını bir dalın ucuna eklememize ya da !!{signed formula}s arasından
ikisini yan yana yerleştirmemize izin verir. İkinci durumda dal ikiye ayrılır
ve eklenen !!{signed formula}s iki yeni dalın ucunu oluşturur.

Karmaşık bir !!{signed formula} için bağlaç kurallarından biri uygulandığında,
!!{signed formula}s biçimindeki dolaysız alt-!!{formula}s eklenir
\iftag{FOL}{; niceleyici kuralları ise uygun yerine koyma örnekleri ekler.}{.}
Kurallar !!{signed formula}s
için, iki işaretin her birine birer tane olmak üzere çiftler hâlinde gelir. Örneğin
$\TRule{\True}{\land}$ kuralı $\sFmla{\True}{!A \land !B}$ ifadesine
uygulanır ve $\sFmla{\True}{!A \land !B}$ ifadesini içeren herhangi bir dalın ucuna
$\sFmla{\True}{!A}$ ile~$\sFmla{\True}{!B}$'yi eklememize izin
verir.
$\TRule{\False}{\land}$ kuralıysa $\sFmla{\False}{!A}$ ile
$\sFmla{\False}{!B}$'yi yan yana ekleyerek bir dalı ikiye ayırmamıza izin
verir. Bir !!^a{tableau}, her dalında !!{signed formula}s arasından eşleşen
$\sFmla{\True}{!A}$ ve $\sFmla{\False}{!A}$ çifti bulunuyorsa kapalıdır.

!!{tableau}s esas alınarak $\Proves$ bağıntısı şöyle tanımlanır:
$\Gamma_0 \subseteq \Gamma$ olacak biçimde sonlu bir
$\Gamma_0 = \{!B_1, \dots, !B_n\}$ kümesi ve şu varsayımlar için kapalı
!!{tableau} varsa $\Gamma
\Proves !A$ olur:
\[
\{\sFmla{\False}{!A}, \sFmla{\True}{!B_1}, \dots, \sFmla{\True}{!B_n}\}.
\]
Örneğin aşağıdaki kapalı !!{tableau}, $\Proves (!A \land !B) \lif !A$
olduğunu gösterir:
\begin{oltableau}
  [\sFmla{\False}{(\formula{A} \land \formula{B}) \lif \formula{A}}, just = \TAss
    [\sFmla{\True}{\formula{A} \land \formula{B}}, just = {\TRule{\False}{\lif}[1]}
      [\sFmla{\False}{\formula{A}}, just = {\TRule{\False}{\lif}[1]}
        [\sFmla{\True}{\formula{A}}, just = {\TRule{\True}{\land}[2]}
          [\sFmla{\True}{\formula{B}}, just = {\TRule{\True}{\land}[2]}, close]
        ]
      ]
    ]
  ]
\end{oltableau}

$\Gamma$ kümesi !!{tableau} hesabında, her $!B_i\in\Gamma$ olacak biçimde
şu varsayımlar için kapalı !!{tableau} varsa tutarsızdır:
\[
\{\sFmla{\True}{!B_1}, \dots, \sFmla{\True}{!B_n}\}.
\]

!!^{tableau}s 1950'lerde Evert Beth ve Jaakko Hintikka tarafından birbirinden
bağımsız olarak geliştirildi; Raymond Smullyan tarafından sadeleştirilip
yaygınlaştırıldı. Bir !!a{tableau} kurmak son derece sistematik bir işlem
olduğundan !!{tableau}s kullanımı çok kolaydır. Sistematik yapıları sayesinde
bilgisayarda uygulanmaya da elverişlidirler. Ne var ki !!a{tableau} çoğu
zaman zor okunur ve ispatlarla bağlantısını görmek kolay olmayabilir. Bu
yaklaşım oldukça geneldir ve birçok farklı mantık için !!{tableau} sistemi
vardır. !!^{tableau}s \iftag{FOL}{!!{structure}s}{değerlemeler} bulmamıza da
yardım eder; bunlar verilen kümelerdeki !!{sentence}s sağlar. Küme
sağlanabilirse kapalı bir !!{tableau} bulunmaz; başka bir deyişle her
!!{tableau} açık bir dal içerir. O dalda uygulanabilecek bütün kurallar
uygulanmışsa \iftag{FOL}{sağlayan !!{structure} açık daldan
“okunabilir”}{sağlayan değerleme açık daldan “okunabilir”}. Sekant hesabıyla
da çok yakın bir bağlantı vardır: kapalı bir !!{tableau}, özünde sekant
hesabında baş aşağı yazılmış ve sıkıştırılmış bir !!{derivation} biçimidir.

